% ECONOMICS_PROBLEM: {"id": "D73-629", "title": "Measuring and reducing corruption", "jel_code": "D73", "source_id": "OP-0125"}
% !TeX program = xelatex
\documentclass[11pt,a4paper]{article}
\usepackage[margin=23mm]{geometry}
\usepackage{fontspec}
\setmainfont{Latin Modern Roman}
\usepackage{amsmath,amssymb,mathtools}
\usepackage{xcolor,enumitem,booktabs,longtable,array,xurl,hyperref}
\definecolor{muted}{HTML}{53636C}
\hypersetup{colorlinks=true,urlcolor=blue,linkcolor=blue}
\setlength{\parindent}{0pt}
\setlength{\parskip}{5pt}
\newcommand{\R}{\mathbb R}
\newcommand{\E}{\mathbb E}
\newcommand{\N}{\mathbb N}
\newcommand{\ind}{\mathbf 1}
\newcommand{\Var}{\operatorname{Var}}
\newcommand{\Cov}{\operatorname{Cov}}
\newcommand{\supp}{\operatorname{supp}}
\DeclareMathOperator*{\argmin}{arg\,min}
\DeclareMathOperator*{\argmax}{arg\,max}
\newcommand{\entrymeta}[3]{{\sffamily\small\color{muted}\textbf{\texttt{#1}}\quad|\quad #2\par #3\par}}
\newcommand{\Problemref}[1]{\texttt{#1}}
\newcommand{\JELref}[2]{\texttt{#2} (JEL \texttt{#1})}
\begin{document}
\section*{Measuring and reducing corruption}
\textbf{Problem ID:} \texttt{D73-629}\quad \textbf{JEL:} \texttt{D73}\par
% BEGIN ECONOMICS STATEMENT
\label{problem:OP-0125}
\label{atlas:prob:Q5}
{\sffamily\small\color{muted}\textbf{Q5} | Political economy | Editorial research problem family\par}
\subsection*{Definitions and assumptions}
Fix public projects $j$, a horizon $h$, and an exhaustive declared set of corruption channels $\mathcal K$. Let $C_{jkh}(a)\geq0$ be net public resources diverted through channel $k$ under monitoring, disclosure or incentive policy $a$. Avoid double counting a transfer moving through several intermediaries. Let $C_{jh}(a)=\sum_{k\in\mathcal K}C_{jkh}(a)$ be integrable true diversion. Reported cases $R_{jh}$ are measurements, not the latent outcome. For a binary occurrence $I_{jh}$, a simple measurement model is
\[
\Pr(R_{jh}=1\mid a)=s_a\Pr(I_{jh}=1\mid a)+f_a\Pr(I_{jh}=0\mid a),
\]
where $s_a$ and $f_a$ are detection sensitivity and false-positive probability; both may change with policy.
\subsection*{Formal research question}
Identify or bound
\[
\theta_C(a,a_0)=\mathbb E[C_{jh}(a)-C_{jh}(a_0)],\quad
\theta_k(a,a_0)=\mathbb E[C_{jkh}(a)-C_{jkh}(a_0)],
\]
jointly with changes in monitoring costs and service quality. Determine which independent audit measurements, assignment designs and restrictions on $s_a,f_a$ make total diversion distinguishable from improved detection or displacement. If choosing a policy, minimize expected diversion plus declared implementation costs over a feasible set; a monetary welfare conversion for service-quality effects must be specified.
\subsection*{Interpretation and scope}
More detected cases may coexist with less true corruption. A reduction in one procurement margin may coincide with increased patronage or diversion elsewhere. Consequently, the set of channels and the observation process are substantive parts of the problem. Informal transfers are not automatically corruption: the legal and economic definition must be fixed. A welfare evaluation must include both recovered resources and any effects on legitimate public-service delivery.
\subsection*{Source audit and status}
All three original citations are verified. Olken measures missing project expenditure using independent engineering assessments; Ferraz--Finan concerns electoral consequences of disclosed audits; Fisman--Miguel studies diplomatic parking behavior. These are distinct outcomes and settings. The total-channel target and measurement model are editorial, retained to correct the atlas's possible conflation of incidence, detection and displacement. Existing results are not relabelled as universally unsolved.
\subsection*{References}
{\footnotesize\raggedright
\begin{enumerate}[label={\textbf{[S\arabic*]}},leftmargin=3em]
\item Benjamin A. Olken (2007). \emph{Monitoring Corruption: Evidence from a Field Experiment in Indonesia}. Journal of Political Economy 115(2), 200--249.\par
\textit{Locator and source role:} Abstract and randomized audits; missing expenditure measured independently in road projects.\par
\url{https://www.journals.uchicago.edu/doi/10.1086/517935}
\item Claudio Ferraz and Frederico Finan (2008). \emph{Exposing Corrupt Politicians: The Effects of Brazil's Publicly Released Audits on Electoral Outcomes}. Quarterly Journal of Economics 123(2), 703--745.\par
\textit{Locator and source role:} Abstract and audit-release timing; electoral accountability conditional on measured corruption, not proof that all corruption channels decline.\par
\url{https://academic.oup.com/qje/article-abstract/123/2/703/1930865}
\item Raymond Fisman and Edward Miguel (2007). \emph{Corruption, Norms, and Legal Enforcement: Evidence from Diplomatic Parking Tickets}. Journal of Political Economy 115(6), 1020--1048.\par
\textit{Locator and source role:} Abstract and diplomatic-parking enforcement setting; evidence on norms and enforcement, with a limited behavioral outcome.\par
\url{https://www.journals.uchicago.edu/doi/10.1086/527495}
\end{enumerate}
}

\subsection*{Common conventions: Source mathematical conventions}

Unless an entry states otherwise, agent and firm sets are finite. State and action spaces are measurable, strategies and outcome maps are measurable, probability laws are specified, and expectations exist for the targets under discussion. Infinite-horizon discount factors lie in $(0,1)$ and discounted objectives are finite. Norms, tolerances and complexity input sizes must be specified for computational comparisons. Constants or primitives that appear locally are inputs, not empirical facts asserted by this catalogue.

\textbf{Identification.} A statistical model $m\in\mathcal M$ implies an observable law $P_m$ and target $\theta(m)$. At law $P$, its identified set is
\[
\Theta(P;\mathcal M)=\{\theta(m):m\in\mathcal M,\ P_m=P\}.
\]
Point identification holds when this set is a singleton, or equivalently when $P_{m_1}=P_{m_2}$ implies $\theta(m_1)=\theta(m_2)$ for all compatible models. Sharp bounds mean bounds attained, or approached when the boundary is not attained, by compatible models. Writing this set is a definition, not a solution: the substantive task is to establish informative restrictions and derive or estimate the set. An unrestricted-confounding model may give only trivial bounds.

\textbf{Inference.} Identification, estimability and valid uncertainty quantification are separate questions. Fix $\alpha\in(0,1)$. A confidence procedure $C_n$ over a declared law class $\mathcal P$ has asymptotic uniform coverage at level $1-\alpha$ when
\[
\liminf_{n\to\infty}\inf_{P\in\mathcal P}\Pr_P\{\theta(P)\in C_n\}\geq1-\alpha.
\]
For set-identified targets the coverage event must be defined separately, for example coverage of the full identified set. If an entry uses identification-indexed models instead of uniquely indexed laws, the same coverage convention is applied over those models. Pointwise limits do not establish uniform validity, and asymptotic validity does not guarantee finite-sample precision. Sample sizes, dependence structures and weak-identification sequences are part of each application.

\textbf{Counterfactuals.} $Y_i(a)$ is a potential outcome under a specified intervention, including its timing, implementation and versions. It is not $Y_i$ conditional on voluntarily choosing $a$. A full intervention vector permits interference; shorthand scalar treatment notation does not silently impose no interference. Assignment, selection, attrition, anticipation and measurement must be specified. A claim about an instrument's local effect is not automatically a population policy effect.

\textbf{Economic equilibrium.} A counterfactual model includes best responses, constraints, feasibility, market clearing, state transitions and expectations consistent with the stated information. When several equilibria exist, either a declared selector is an input or the target is set-valued. Comparisons across policy regimes must use the stated selection convention. Reduced-form relationships are not presumed invariant after an equilibrium-changing intervention.

\textbf{Optimization and welfare.} Optimization is over the stated feasible set. If attainment is not established, $\sup$ or $\inf$ and $\varepsilon$-optimal choices replace an asserted maximizer or minimizer. For example, with value $V=\sup_{a\in\mathcal A}W(a)<\infty$, an $\varepsilon$-optimal set is $\{a\in\mathcal A:W(a)\geq V-\varepsilon\}$ for $\varepsilon>0$. Benefits, costs, risk and health or environmental outcomes can be aggregated only using declared units and conversion weights. Interpersonal utility weights, the treatment of fiscal transfers and foreign profits, discounting, and the behavioral welfare criterion are explicit inputs. A comparative-static derivative additionally requires existence and appropriate differentiability of the selected counterfactual.

\textbf{Sources.} Local labels [S1], [S2], and so forth restart in every question. Locators identify the relevant abstract, model, section or result at the level actually checked; a bibliographic or abstract-level verification is not represented as inspection of an entire proof. Working papers are distinguished from later publications and changing versions. Source summaries are paraphrased. All URL checks and editorial formulations refer to the audit date stated above.
% END ECONOMICS STATEMENT
\end{document}
